\section{The Twelfth and Final Labor}

We have finally arrived at the last leg of the journey in the Dodekathalon. Hercules is given a final test, as he ``cheated'' on two of the earlier ones, and this one is meant, not to test his mettle, but to ruin and destroy him. There is a chance, if he truly goes to Hades, that he may not return, as even the godling Herakles is one of the ``living'' who may not be allowed to enter, and then just waltz out of hell. Hell is meant to keep the dead and the living separate; you don't keep an outhouse inside the bedroom, and you don't allow the living and the dead worlds to become indiscriminately intermingled.

In the ancient world, such sages as Pythagoras or Socrates were said to have been accompanied and guided by their daemon, who instructed them in their journey toward the truth. Of a truth, we could say that Hercules has been guided by a kind of warrior-appropriate daemon, in that he wars through his passions (as Guido De Giorgio\footnote{\url{https://gornahoor.net/?p=6036}} has written) ``on his own'', without the conscious guide of a familiar spirit or guardian angel. The warrior has to perform the self-guiding function on his own by conquering passions and his body. We note\footnote{\url{https://en.wikipedia.org/wiki/Labours\_of\_Hercules\#Twelfth\_Labour:\_Cerberus}}, however, that even the warrior has occasional need of the mysteries: Hercules goes for instruction to Eleusis.

It is said by Iamblichus (in his biography\footnote{\url{http://archive.org/stream/IamblichusPythagoreanLifetaylor/Iamblichus-PythagoreanLife\_djvu.txt}} of Pythagoras) that he imposed the most careful methods of scrutiny and tests on those of his auditors who were judged progressively worthy to receive his teachings. He received his teachings from Egypt, a land which was under the dominion of Mercury, but he passed into Italy, as this was a land renowned for its upright morality (Pythagoras was ``drawn'' by affinity towards the simple, basic, and rare quality of morality). There is a link, here, too, to Hercules, because he came to Crotona\footnote{\url{http://books.google.com/books?id=pavWAAAAMAAJ&pg=PA15&lpg=PA15&dq=crotona+hercules+labors&source=bl&ots=sEfUAs4J9f&sig=dNeWlkpPOgK\_JyXdGEsllzAIYZM&hl=en&sa=X&ei=4XbiUtTcLIuhsQS9k4H4DQ&ved=0CDoQ6AEwAw\#v=onepage&q=crotona\%20hercules\%20labors&f=false}} and taught them, a city founded by Hercules during his travels. No one was treated ``equally'' by Pythagoras -- although most had to endure apprenticeships and tests, Abaris\footnote{\url{https://en.wikipedia.org/wiki/Abaris\_the\_Hyperborean}} the Hyperborean who came to the city was judged by him immediately worthy of his highest discourses, and they exchanged knowledge and gifts, Abaris giving him a dart with which Pythagoras threw out the diseases from Laecadomonia, and Pythagoras showing him his golden thigh, proof that he was attended by Apollo and was a hero or demi-god.

From such caste of brahmin, at Eleusis (or those trained or descended spiritually from them), Hercules can hope to learn much. Indeed, he finds an entrance into the underworld, as well as advice about how to safely return. But here, in the Underworld, Hercules is first and foremost a suppliant. He must take and receive permission from Hades to attempt to wrestle Cerebus back to the surface. The warrior, even the demi-god, doesn't exalt himself in the presence of the ruler of the Dead, whose task is inexorable, but rather consults with him and presents a petition or craves a boon. Hades grants the wish, on condition that no weapon is used.

Incidentally, Hercules rescues one hero from hell, \& fails to rescue the other, whose desire to have a god's wife for his own was so insulting that nothing could transport the man out of hell. The gravity of karma in the afterlife has certain laws that even a demigod can't bend. This is why Scripture tells us not to fear those who rend the body, but the One Who can chain you in Hell. Don't be so heavy in the afterlife that even a god can't haul you out of prison. All of this indicates that the laws of Heaven and Hell (and, in between, Middle Earth) are more important and complex than even some of the tales of the demigods. Hercules' story itself is a foreshadowing of greater deeds to come, when Hell will be broken open and harrowed in truth. This does not detract from Hercules, but rather adds luster and glory to the daring deed, which might otherwise be without absolute significance.

Naturally, Hercules brings down the big dog Cerebus, and wrests him from his hellish lair up to the sweet light of the sun. Just as inevitably, Eurystheus flees back into this pithos at the sight of this still unexpected victory. Can we begin to suspect, although the legends do not say so, that Hercules (like Socrates) possessed a daemon that guided him in his wrestling and fighting, instructing him minutely and instantly upon what to do? In any case, he answers to a ``higher power'' either within or partially without him.

Alice Bailey has some reasonable thoughts\footnote{\url{http://www.souledout.org/hercules/herculescapricorn/herccap.html}} on Capricorn and the twelfth labor, specifically the idea that the initiate has to go through a personal hell before he can instruct others (which involves traversing objective Hell). If you can't be pure in your own circle, you cannot purify the world either. This is a strong argument for the emphasis on Law which Gornahoor has made. That is (for example) why would America ``save'' the world when her own streets are in chaos and disorder? We have seen that Boris Mouravieff in Gnosis implicitly endorses Order over Chaos as a necessary prerequisite to any gnosis itself, when he stated that in some conditions politically the chaos becomes too great for any ``work''. St. Paul endorses Law in Romans 13 and other passages --- pray for a peaceable life of quiet godliness. Pythagoras endorsed the Romanity of the Latins by emigrating to Crotona from turbulent Samos, when the revolution arrived, and he taught his initiates to favor those of Law and to be unfriendly to those who advocated Revolution. Or, as the German folk used to say, work first, then play or relaxation.

A lot of the New Age followers, or modern people interested in initiation are short circuiting the process and cheating themselves (not to mention defiling the memory of their ancestors) by holding to political doctrines that are simply incompatible with ``decency and good order''. Far better to be apolitical in the fashion of Ernst Junger, in order to achieve interiority, or better still, give your reason to the meditation of Order. Even Hercules did this, in coming hat in hand to Hades, in consulting the Eleusian brahmins/priests. Christ Himself, God Himself, did not consider Equality a thing to be grasped! But rather humbled himself as a servant, and said that not one jot or tittle of the General Law would be cheated, but rather that He intended to sacrifice His flesh and soul and humble His divine origin in order to carve a way out through His own flesh and interiority. It is said that Satan does not hate man, but simply has no faith in human nature, and that this is why he ``opposes'' arrogant man, tempting him. ``Every man is tempted out of his own flesh''. The deities and powers that exist take note of those, especially in a dark time, who observe the remnants of the Law and do not bow the knee to man-made Baal. The Creator, the Archetype, is not mocked, especially by those who think that they can take a short cut away from Law in favor of Chaos, which is Non-Being. Sow the wind, reap the whirlwind. Or, like Gandalf, place your faith in the little things and aid those who bear the ring faithfully.

The Pythagoreans believed that Number was the motion of the Gods, or perhaps the essence of the Gods themselves (Number as the Intelligible: number as sensible would be the tracing of their steps), so their philosophy excluded Chaos and Revolution as non-Being, since Being was orderly, numerate. Here we see that all initiates of Truth agree in rejecting anti-Order; indeed, a great Reformation thinker like van Prinsterer demonstrated\footnote{\url{http://www.reformationalpublishingproject.com/pdf\_books/Scanned\_Books\_PDF/UnbeliefinReligionandPoliticsUnbeliefandRevolutionLectures\_VIII\_IX.pdf}} that the Revolution was anything but natural, inevitable, or (therefore) ``orderly'' --- rather, it was an attempt to (as Pythagoras would have said) ignore the Creator and pay attention to sublunary magistrates in the place of the Creator. This is unseemly, with predictable eventual results.

If Christ opened the path to the heavens, \emph{recollecting} the content of this gift, in meditating the ancient mysteries and the old paths that are to be walked in, is the surest way available to understand a gift which otherwise is easily misinterpreted as just another gigantic Revolution or shortcut to cheat the Lord of Lords. Christ did not come to create a new world, but to re-create the Old.

I don't endorse everything \emph{Civilizing the Beast}\footnote{\url{http://civilizingthebeast.blogspot.com/2013/06/ordered-evolution-rather-than-ordered.html}} writes, but he is right to seek a way out of the Left's empty insistence that ``man as he is'' (or is understood right now) can be enshrined in a Constitution as a secular demi-god, or the new Right's unintentional reversion to exoteric forms of the Past which no longer have spirit or life left within them. If Christ is holding exoteric/esoteric worlds together with the cross, this means that the Dreamtime\footnote{\url{https://en.wikipedia.org/wiki/Dreamtime}} will become the wide awake truth of the Future, which always was to begin with, through the Man. It means that man will become God, as the Fathers taught. It means that history is becoming more intricately complex and subtle every day, as the Spirit weaves where it wills. Instituting the ruins of man and his revolution as ``Law'' is the besetting hubris of classical liberalism, since man is not merely his discursive Reason.

If you want to get to heaven, you've got to raise a little hell\footnote{\url{https://www.youtube.com/watch?v=kYvOsnhV6ZY&list=PL5DACA82DCA49E11B}}. Start, and end, with yourself. Run the race, as if to win. But do not neglect the ``little things'', as each adventure of the twelve labors encapsulates the whole, and even the mighty Hercules began with a ``little labor''.


\flrightit{Posted on 2014-01-25 by Logres}

\begin{center}* * *\end{center}

\begin{footnotesize}
\begin{sffamily}
\texttt{Paulo on 2014-01-26 at 02:15 said:}

Were are the men to be found to accomplish this task?They will be seen when the separation of the wheat the chaff start to occurs?Those who have eyes will see!

\hfill

\texttt{Teuton on 2014-02-01 at 23:03 said:}

This isn't directly related to this post, but I wonder, is there any means by which readers could contact the authors of this blog to discuss various issues or questions they might have regarding various aspects of Tradition or the content of the posts here? I know there used to be a forum here, but that seems gone.

I wonder simply because I have questions that I think authors here might be able to answer, and I see this website as producing higher quality and more informed articles and insights than one finds elsewhere. In other words, I respect the content here.

Just for an example, I have some questions about the concept of initiation and how necessary it is for a valid spiritual path for someone who seeks gnosis. It seems Guenon and others highly stress its importance, but what of people who have had authentic experience of transcendence? If someone, by God's grace, had an authentic transcendent mystical experience, is such an experience equivalent to or akin to initiation? Thus having experienced a ``glimpse'' of Truth or the ultimate goal, could one proceed along a path that might not have an esoteric initiatory chain, say like traditional Roman Catholicism (I am still unsure if the sacraments are initiatory in the Guenonian sense or not.) I wonder because I am not really interested in losing my Western heritage by going East, even as far east as the Eastern Orthodox world, but like all true seekers I wish to find the Way. It seems modern Catholicism has a plethora of problems (but then which tradition doesn't in our dark age?), likely necessitating a Sedevacantist positioning in regards to the modern Popes, exoteric though that concern may be. Then again, it seems that, degenerated though Catholicism is, it is a valid path since it produced figures like Eckhart and Ruysbroeck historically, and even if by pursuing these lofty aims one might be a complete minority or lone wolf in relation to ones brothers in faith.

Thanks.

\hfill

\texttt{Logres on 2014-02-04 at 12:15 said:}

Good questions, Teuton. I tend to agree that even Orthodox is too far East for my soul. Celtic Orthodoxy has a certain amount of interest as well.

\hfill

\texttt{Teuton on 2014-02-08 at 12:27 said:}

Well, perhaps either you or Cologero could either comment on or point me in the direction of a better understanding of initiation. Again, I wonder if someone actually had a true experience of transcendence/verticality, an experience of timelessness/spacelessness and of a supra-personal reality beyond all concepts, etc. whether or not such an experience would be the equivalent to initiation. If such a one were unable to transmit that experience to another initiate, I'd presume not, but one wonders if such an experience would function as a valid ``seed'' toward the fruit of gnosis and how it ultimate would relate to initiation if, again, it were a true transcendent experience.

\hfill

\texttt{Logres on 2014-02-09 at 00:50 said:}

Teuton, I can't speak for Cologero, although he has addressed this question in many and several various places, both in comments and essays. There are several issues at stake.

1. People have experiences all the time -- because they aren't prepared, they don't go far enough, and they often cannot interpret them. So the real issue is preparation, whether this life or after death (post mortem). I think Cologero has said this, and BTW, all of what I know I owe to him, on these subjects, indirectly at the least.

2. The initiation from above, especially in our age (a spiritual winter, or Kali Yuga, when the sun withdraws from the earth, and things coagulate in the cold), is often transmitted in unlooked for places. Think of winter, with Christmas at its heart.

3. Mystical experience and initiation are qualitatively different, so Cologero may point out other things that apply specifically to initiation.

4. There is big debate over whether Christian sacraments are a form of initiation: Guenon says no, Schuon said yes -- a French order was ``put into sleep'' over the debate, so as not to induce schism. Cologero has commented on some weaknesses in Guenon, but doesn't comment much on Schuon. Evola has an even different take.

You are doing the right thing, asking the right question. Just keep sharpening it. You can't get the right answers to the wrong questions. I recommend Cologero's comments recently on Guenon, over the last several months, as very enlightening, as they represent an adjunct or commentary (final thoughts) on someone he regards as the ``Master of Tradition''. I hope this helps.

\hfill

\texttt{X on 2014-02-09 at 07:59 said:}

``Mystical experience and initiation are qualitatively different''

Logres, can you please define the difference between ``mystical experience'' and ``initiation'' here? Did you mean the former as the synonym of Mysticism?

\hfill

\texttt{Logres on 2014-02-09 at 10:24 said:}

That's right.

\hfill

\texttt{Teuton on 2014-02-09 at 12:37 said:}

Thanks for the reply. I ask because I had a spontaneous experience around 7 years ago which changed me from an ambivalent atheist (someone who just didn't care for these subjects) to a spiritual seeker with gnosis as my ultimate goal in life overnight. I didn't know what happened to me, but I did know that it was sacred and qualitatively different from anything one can experience in the mundane world of day-to-day experience, i.e. that it was certainly transcendent in the real sense.

In order to discover what happened to me I began to research various topics, and once I read the Upanishads I felt the sages were discussing the reality I had a ``glimpse'' of, whereas for them it was a permanent attainment. Later I discovered Evola and most of what he said (aside from his overly critical attitude toward Christianity in my opinion) struck a deep chord within me.

My only issue, as described earlier, is that I don't hate my race, ethnicity, or heritage and would like to retain my Western identity. Beyond that, I am skeptical about how healthy the Eastern traditions are in any case. My concern was, having an experiential reference point, whether initiation is an absolute requirement to become a ``jivanmukta'' as the Hindus say, since it seems it is often portrayed that way.

That said, I will certainly look over the material by Cologero you suggested.

\hfill

\texttt{Logres on 2014-02-09 at 17:51 said:}

\url{http://www.gornahoor.net/?p=5578}

I need to reread this one, by Mercurius.
\url{http://www.gornahoor.net/?p=5116}

Cologero put this up, by Coomaraswamy.
\url{http://www.gornahoor.net/?p=5605}

On preparation being paramount... 
\url{http://www.gornahoor.net/?p=6930}

Example of his meditations on Guenon

I think there are more, and I seem to remember one being particularly good I can't find right now. You are doing the first step: trying to remember what your experience was, by clarifying it, without vitiating it, so that you can make the best use of it. It's a wise start! Glad to hear people still have such experiences. While I was in college, during a particular low point, while walking and praying, and having some agnostic intellectual problems that were devouring my faith, I saw something come out of a graveyard that convinced me that the invisible world was very tangible and real. After that, I never had any intellectual doubts about God's existence. But had I not been reading philosophy, I might have taken a very New Age turn.

\hfill

\texttt{scardanelli on 2014-02-11 at 08:19 said:}

Teuton,

I came across this and thought it might be helpful for you. My understanding, and that which I take as truth, is the initiation of the Second Birth in which Christ is the only initiator. This initiation happens individually and not necessarily as a part of an established group or order. One find's this idea expressed in a certain ``school'' of writers, notably Tomberg and many of those mentioned by him in MOTT. In ``History and Doctrines of the Rose-Croix'' Sedir states that :

``Nowadays, whatever their leaders say, secret societies or perhaps orders, initiatory organizations, etc are in their dotage, at least in our country; and people in general are slowly being transformed in their collective organisms and little by little, are becoming capable of establishing communication with the invisible in their consciousness, in broad daylight. These developments are destined to grow incessantly, until the blessed dawn when the name of the Father shall be hallowed on Earth as it is in Heaven.''

As Logres suggested, read through the archives here and at \url{http://www.meditationsonthetarot.com}. If you look carefully you will find a golden thread woven throughout. It is perhaps better to find on your own than to simply have all of your questions answered.

\hfill

\texttt{Aleksandar Gueric on 2015-10-12 at 12:41 said:}

Great article with so many elucidating insights, I will have to re-read it few more times. This one and few others written by Cologero helped me a lot to crystallize thoughts and ideas that were haunting me for years. Namely, how to harmonize my Christian and certain pre Christian parts of my belief. Answer seems so simple now -- Traditional Christianity is not just a continuation of one and true Tradition, but it is a religion of our Age, Kali yuga or Wolf age.

I've always felt that Christianity represents a crystalization of pagan tradition, of entire human history and existence up to that point. It's the last myth, Christ on Earth, Logos incarnated and descended to Middle Earth to prepare and transfer humanity into a new and last age of our cycle. Who would know it would be so dark?

Maybe Heracles was a Christ of his Age? Son of God, Solar principle and '\,'miracle worker'\,' who died only to reach apotheosis? As he paved a path of Olympians through pure virility, out of the decadent chaotic matriarchy, humanity progressed and become ready for another, final transformation in the midst of the darkest night. This is a Hyperborean path, ultimately. And speaking of Hyperborea, I'm glad you also mentioned Abaris. I always wanted to know about this mysterious figure and priest of Apollo.

Both Apollo and Heracles, while '\,'young'\,', killed a serpent. I'm fascinated by this motif of a solar warrior slaying a chthonic serpent/dragon. When I was very young, an old hag living next door forced me to kill a snake in her garden. I used a pitchfork she gave me. I remember regreting it very much , because everyone told me it's bad to kill a snake. I relate that now to a symbolism of solar vs chthonic struggle. It's good to mythologize your own childhood sometimes. Or enchant our life and world around. Still, images of chthonic nature surf out of my subconsciousness and reveal themselves for a moment. Images of titanic serpents trying to swallow me, abysmal caverns and disturbing saurian and insectoid presence. I interpret it as it's just my subconsciousness emerging to the surface. I think that often my subconsciouss comes dangerously near the surface. Especially when i'm in the contact with water. No wonder I drowned 3 times although I'm a good swimmer.

On the other hand, Solar Olympian imagery induces a trance like state and I feel elevated in spirit. Reading Evola has an erotic character for me. Same goes for Plato or even Novalis. I remember reading *Revolt Against the Modern world* for the first time, vivid visions of columns and Olympian temples and symmetric, archetypal Aryan architecture on snowy high mountain peaks danced in front of me as I felt ecstatic. It is my metaphysical exile, as Cioran would say.

If this life is a test, and it has to be, I fear that most people do not even realize it. They keep walking, not even trying to win the race. My mind is constantly aware of it and it feels very scary sometimes. I can not turn it off anymore, seeing myself as an image in an abstract dimension. Are people on the street even aware of life? My life presented me very fine Heraklian labors and I am in front of another big one. With the new strength and knowledge, thanks to this site as well, I feel more prepared for new challenges. Time seems to speed up and chaos is growing each day. Jesus Prayer brings tears no matter how hard I try to stop it. Truly dark times are ahead of us.


\hfill

\end{sffamily}
\end{footnotesize}

